\section{Background and related research}
\label{sec:background}

An ontology describes categories, relationships, and constraints within a
domain. An ontology reasoner determines what follows from these descriptions
and the available observations. We use the thesis's original Bach family
knowledge base as a running illustration~\cite[Table~2.5, p.~35]{volz2004}.
It describes Johann Ambrosius as a man with a male child, without naming that
child. It also states that men are persons and defines a father as a man with
a child who is a person. A reasoner can therefore conclude that Johann
Ambrosius is a father without identifying his child. This distinguishes an
implied conclusion from an explicitly recorded statement.

Missing information is also different from falsehood. The same knowledge base
names Wilhelm Friedemann as Johann Sebastian's child. As the thesis explains,
this does not establish that he is Johann Sebastian's only child
\cite[Example~5.4.1, p.~131]{volz2004}. Incomplete family records do not rule
out additional children.

A complementary illustration comes from the Lehigh University Benchmark
(LUBM), used in our external comparisons~\cite{guo2005,lubmwebsite}.
Its ontology states that every graduate student is a person who takes some
graduate course, that graduate courses are courses, and that people taking
courses are students. A graduate-student declaration therefore implies
student membership even when it names no course. This distinction matters
for LUBM's sixth query, which asks for all students. Bach connects the present
study to the original work; LUBM connects the same reasoning ideas to an
established benchmark.

The Web Ontology Language, OWL, gives such statements a precise meaning
\cite{owl2004,owldirect2012}. Greater expressive power permits more complex
descriptions, but can make inference more demanding. A central research
question is therefore which restrictions preserve useful forms of expression
while enabling simpler reasoning. A \emph{fragment} is a language obtained by
such restrictions. Its value can arise from a mathematical guarantee, a
computational advantage on a specified task, or both.

Description Logic Programs were introduced as a connection between ontology
languages and rules of the form ``if these conditions hold, then this
conclusion follows''~\cite{grosof2003}. Volz's 2004 thesis developed this
connection into methods for answering questions and maintaining inferred
knowledge as facts and rules change~\cite{volz2004}. The original research
already addressed changes to the ontology's rules, rather than only changes
to observations~\cite{volz2003maintenance}. The work examined here inherits
that problem and semantic foundation. Its expressive mode, called DLP L3,
also admits statements requiring an individual to exist without naming them.
The thesis illustrates this using Johann Ambrosius's unnamed child
\cite[Example~4.5.2, p.~94]{volz2004}. This additional capability requires separate
attention to whether a reasoning procedure terminates.

One influential approach computes consequences in advance and retains them
for later questions. This is called \emph{materialization}. Its main attraction
is that repeated questions can reuse earlier deductions. Its main difficulty
is maintaining those deductions when their premises change. Adding a new
premise usually creates further conclusions. Removing one is subtler: a
conclusion may have several independent justifications. The thesis's extended
Bach family tree supplies two paths from Johannes to Wilhelm Friedemann, one
through Johann Sebastian and another through Maria Barbara. In its hypothetical
withdrawal of the Johann Sebastian--Wilhelm Friedemann link, Johannes remains
an inferred ancestor through Maria Barbara
\cite[Section~6.2.1 and Figure~6.2, p.~150]{volz2004}. The example changes the
recorded knowledge, rather than making a historical claim about parentage.

Research has developed this idea in several directions. Deletion and
rederivation, established before the thesis, first withdraw potentially affected
conclusions and then restore those still justified~\cite{gupta1993}.
Subsequent methods examine alternative proofs more selectively, count suitable
forms of support, or combine specialized procedures
\cite{motik2015bf,hu2018,motik2019,hu2022modular}. Changes to the rules
themselves have also received renewed attention
\cite{kotowski2011,xu2023}. These studies establish the broader principle
that maintenance should depend on the structure of the affected reasoning,
not simply on the number of changed statements.

A second direction reduces repeated work across similar deductions.
Type-based reasoning groups individuals with recurring descriptions, while
compressed reasoning represents repeated patterns more economically
\cite{bajraktari2017,hu2019compressed}. Parallel and column-oriented
systems organize computation and storage to handle larger knowledge bases
\cite{motik2014parallel,carral2019}. More recent work further specializes
representations for common patterns~\cite{zhang2024}. These approaches
motivate our investigation of whether individuals with the same current
starting facts can share the work of establishing a surviving justification.
The general idea of sharing proofs is inherited; the present study tests a
restricted form of it during simultaneous changes to facts and rules.

A third direction concerns the order in which conditions are tested. Combining
records that satisfy several conditions is known as a \emph{join}. An
unfavourable order can produce many temporary combinations that later tests
discard. Research has established stronger algorithms for combining multiple
relations~\cite{ngo2012}, and recent studies investigate robust filtering
and the avoidance of unnecessary intermediate expansion
\cite{qiao2026,bekkers2025}. These findings motivate our experiments with
limited adaptation of evaluation order. Their mathematical guarantees do not
automatically apply to the simpler strategy examined here.

Finally, OWL 2 profiles formalize different restrictions for different
reasoning needs. The rule-oriented RL profile explicitly acknowledges DLP
among its influences~\cite{owlprofiles2012}. This historical continuity
motivates reassessment rather than a claim that the earlier fragments have
become obsolete. Our contribution is a corrected reconstruction and a
controlled investigation of when selected restrictions and reasoning
strategies help, including cases in which their costs outweigh their benefits.
